\documentclass{amsart}
% For dealing with references we use the comment environment
\usepackage{verbatim}
\newenvironment{reference}{\comment}{\endcomment}
%\newenvironment{reference}{}{}
\newenvironment{slogan}{\comment}{\endcomment}
\newenvironment{history}{\comment}{\endcomment}

% For commutative diagrams we use Xy-pic
\usepackage[all]{xy}

% We use 2cell for 2-commutative diagrams.
\xyoption{2cell}
\UseAllTwocells

% We use multicol for the list of chapters between chapters
\usepackage{multicol}

% This is generally recommended for better output
\usepackage{lmodern}
\usepackage[T1]{fontenc}

% For cross-file-references

% Package for hypertext links:
\usepackage{hyperref}

% For any local file, say "hello.tex" you want to link to please

% Theorem environments.
%
\theoremstyle{plain}
\newtheorem{theorem}[subsection]{Theorem}
\newtheorem{proposition}[subsection]{Proposition}
\newtheorem{lemma}[subsection]{Lemma}

\theoremstyle{definition}
\newtheorem{definition}[subsection]{Definition}
\newtheorem{example}[subsection]{Example}
\newtheorem{exercise}[subsection]{Exercise}
\newtheorem{situation}[subsection]{Situation}

\theoremstyle{remark}
\newtheorem{remark}[subsection]{Remark}
\newtheorem{remarks}[subsection]{Remarks}

\numberwithin{equation}{subsection}

% Macros
%
\def\lim{\mathop{\mathrm{lim}}\nolimits}
\def\colim{\mathop{\mathrm{colim}}\nolimits}
\def\Spec{\mathop{\mathrm{Spec}}}
\def\Hom{\mathop{\mathrm{Hom}}\nolimits}
\def\Ext{\mathop{\mathrm{Ext}}\nolimits}
\def\SheafHom{\mathop{\mathcal{H}\!\mathit{om}}\nolimits}
\def\SheafExt{\mathop{\mathcal{E}\!\mathit{xt}}\nolimits}
\def\Sch{\mathit{Sch}}
\def\Mor{\mathop{\mathrm{Mor}}\nolimits}
\def\Ob{\mathop{\mathrm{Ob}}\nolimits}
\def\Sh{\mathop{\mathit{Sh}}\nolimits}
\def\NL{\mathop{N\!L}\nolimits}
\def\CH{\mathop{\mathrm{CH}}\nolimits}
\def\proetale{{pro\text{-}\acute{e}tale}}
\def\etale{{\acute{e}tale}}
\def\QCoh{\mathit{QCoh}}
\def\Ker{\mathop{\mathrm{Ker}}}
\def\Im{\mathop{\mathrm{Im}}}
\def\Coker{\mathop{\mathrm{Coker}}}
\def\Coim{\mathop{\mathrm{Coim}}}

% Boxtimes
%
\DeclareMathSymbol{\boxtimes}{\mathbin}{AMSa}{"02}

%
% Macros for moduli stacks/spaces
%
\def\QCohstack{\mathcal{QC}\!\mathit{oh}}
\def\Cohstack{\mathcal{C}\!\mathit{oh}}
\def\Spacesstack{\mathcal{S}\!\mathit{paces}}
\def\Quotfunctor{\mathrm{Quot}}
\def\Hilbfunctor{\mathrm{Hilb}}
\def\Curvesstack{\mathcal{C}\!\mathit{urves}}
\def\Polarizedstack{\mathcal{P}\!\mathit{olarized}}
\def\Complexesstack{\mathcal{C}\!\mathit{omplexes}}
% \Pic is the operator that assigns to X its picard group, usage \Pic(X)
% \Picardstack_{X/B} denotes the Picard stack of X over B
% \Picardfunctor_{X/B} denotes the Picard functor of X over B
\def\Pic{\mathop{\mathrm{Pic}}\nolimits}
\def\Picardstack{\mathcal{P}\!\mathit{ic}}
\def\Picardfunctor{\mathrm{Pic}}
\def\Deformationcategory{\mathcal{D}\!\mathit{ef}}

\setlength{\emergencystretch}{3em}
\begin{document}
\title[Stacks Project, Tag 0C54]{579 --- Completed local presentations extend coefficient fields under dimension-preserving field base change with separable original residue field}
\maketitle
\noindent Copyright (C) 2005--2025 Johan de Jong. Stacks Project authors. Source: \href{https://stacks.math.columbia.edu/tag/0C54}{Tag 0C54}.

\noindent Editorial difficulty note (not author text). Difficulty H2: The proof must construct a compatible coefficient-field embedding after arbitrary field extension. It shows that the local tensor-product fibre is a field, obtains a surjection of complete power-series presentations, and invokes flatness of completion to prove that no additional relations appear..

\noindent Editorial provenance note (not author text). History: excerpt from revision \texttt{a0444\allowbreak 6e57e\allowbreak c1fbc\allowbreak 252a8\allowbreak 71afc\allowbreak ec775\allowbreak 2fb28\allowbreak 07b14}, varieties.tex, lines 4100--4193. Reference presentation was adapted; original-author context and core support blocks were retained or added. The mathematical wording is unchanged. Licensed under GNU FDL 1.2 or later; no invariant sections or cover texts. See ../licenses/COPYING. Context blocks reproduce author definitions and supporting results; they are not additional counted problems.

\begin{lemma}
\label{lemma-complete-local-ring-structure-as-algebra}
Let $k$ be a field. Let $X$ be a locally Noetherian scheme over $k$.
Let $x \in X$ be a point with residue field $\kappa$.
There is an isomorphism
\begin{equation}
\label{equation-complete-local-ring}
\kappa[[x_1, \ldots, x_n]]/I \longrightarrow \mathcal{O}_{X, x}^\wedge
\end{equation}
inducing the identity on residue fields.
In general we cannot choose (\ref{equation-complete-local-ring})
to be a $k$-algebra isomorphism. However, if the extension $\kappa/k$
is separable, then we can choose
(\ref{equation-complete-local-ring}) to be an isomorphism of $k$-algebras.
\end{lemma}

\begin{lemma}
\label{lemma-base-change-complete-local-ring}
Let $K/k$ be an extension of fields. Let $X$ be a locally algebraic
$k$-scheme. Set $Y = X_K$. Let $y \in Y$ be a point with image $x \in X$.
Assume that $\dim(\mathcal{O}_{X, x}) = \dim(\mathcal{O}_{Y, y})$
and that $\kappa(x)/k$ is separable.
Choose an isomorphism
$$
\kappa(x)[[x_1, \ldots, x_n]]/(g_1, \ldots, g_m) \longrightarrow
\mathcal{O}_{X, x}^\wedge
$$
of $k$-algebras as in (\ref{equation-complete-local-ring}).
Then we have an isomorphism
$$
\kappa(y)[[x_1, \ldots, x_n]]/(g_1, \ldots, g_m) \longrightarrow
\mathcal{O}_{Y, y}^\wedge
$$
of $K$-algebras as in (\ref{equation-complete-local-ring}). Here we use
$\kappa(x) \to \kappa(y)$ to view $g_j$ as a power series
over $\kappa(y)$.
\end{lemma}

\begin{proof}
The local ring map $\mathcal{O}_{X, x} \to \mathcal{O}_{Y, y}$
induces a local ring map
$\mathcal{O}_{X, x}^\wedge \to \mathcal{O}_{Y, y}^\wedge$.
The induced map
$$
\kappa(x) \to \kappa(x)[[x_1, \ldots, x_n]]/(g_1, \ldots, g_m)
\to \mathcal{O}_{X, x}^\wedge \to \mathcal{O}_{Y, y}^\wedge
$$
composed with the projection to $\kappa(y)$ is the canonical
homomorphism $\kappa(x) \to \kappa(y)$.
By Lemma \href{https://stacks.math.columbia.edu/tag/0C4Y}{lemma change fields flat (Tag 0C4Y)} the residue field
$\kappa(y)$ is a localization of $\kappa(x) \otimes_k K$
at the kernel $\mathfrak p_0$ of $\kappa(x) \otimes_k K \to \kappa(y)$.
On the other hand, by Lemma \href{https://stacks.math.columbia.edu/tag/0C50}{lemma change fields algebraic unramified (Tag 0C50)}
the local ring $(\kappa(x) \otimes_k K)_{\mathfrak p_0}$
is equal to $\kappa(y)$. Hence the map
$$
\kappa(x) \otimes_k K \to \mathcal{O}_{Y, y}^\wedge
$$
factors canonically through $\kappa(y)$. We obtain a commutative
diagram
$$
\xymatrix{
\kappa(y) \ar[rr] & & \mathcal{O}_{Y, y}^\wedge \\
\kappa(x) \ar[r] \ar[u] &
\kappa(x)[[x_1, \ldots, x_n]]/(g_1, \ldots, g_m) \ar[r] &
\mathcal{O}_{X, x}^\wedge \ar[u]
}
$$
Let $f_i \in \mathfrak m_x^\wedge \subset \mathcal{O}_{X, x}^\wedge$
be the image of $x_i$. Observe that
$\mathfrak  m_x^\wedge = (f_1, \ldots, f_n)$ as the map is surjective.
Consider the map
$$
\kappa(y)[[x_1, \ldots, x_n]] \longrightarrow \mathcal{O}_{Y, y}^\wedge,\quad
x_i \longmapsto f_i
$$
where here $f_i$ really means the image of $f_i$ in $\mathfrak m_y^\wedge$.
Since $\mathfrak m_x \mathcal{O}_{Y, y} = \mathfrak m_y$
by Lemma \href{https://stacks.math.columbia.edu/tag/0C50}{lemma change fields algebraic unramified (Tag 0C50)}
we see that the right hand side is complete with respect to
$(x_1, \ldots, x_n)$ (use Algebra, Lemma
\href{https://stacks.math.columbia.edu/tag/05GG}{lemma hathat finitely generated (Tag 05GG)} to see that
it is a complete local ring).
Since both sides are $(x_1, \ldots, x_n)$-adically complete
our map is surjective by
Algebra, Lemma \href{https://stacks.math.columbia.edu/tag/0315}{lemma completion generalities (Tag 0315)}
as it is surjective modulo $(x_1, \ldots, x_n)$.
Of course the power series $g_1, \ldots, g_m$
are mapped to zero under this map, as they already map to zero
in $\mathcal{O}_{X, x}^\wedge$. Thus we have the commutative diagram
$$
\xymatrix{
\kappa(y)[[x_1, \ldots, x_n]]/(g_1, \ldots, g_m) \ar[r] &
\mathcal{O}_{Y, y}^\wedge \\
\kappa(x)[[x_1, \ldots, x_n]]/(g_1, \ldots, g_m) \ar[r] \ar[u] &
\mathcal{O}_{X, x}^\wedge \ar[u]
}
$$
We still need to show that the top horizontal arrow is an isomorphism.
We already know that it is surjective. We know that
$\mathcal{O}_{X, x} \to \mathcal{O}_{Y, y}$ is flat
(Lemma \href{https://stacks.math.columbia.edu/tag/0C4Y}{lemma change fields flat (Tag 0C4Y)}), which implies that
$\mathcal{O}_{X, x}^\wedge \to \mathcal{O}_{Y, y}^\wedge$ is flat
(More on Algebra, Lemma \href{https://stacks.math.columbia.edu/tag/0C4G}{lemma flat completion (Tag 0C4G)}).
Thus we may apply Algebra, Lemma \href{https://stacks.math.columbia.edu/tag/00ME}{lemma mod injective (Tag 00ME)}
with $R = \kappa(x)[[x_1, \ldots, x_n]]/(g_1, \ldots, g_m)$,
with $S = \kappa(y)[[x_1, \ldots, x_n]]/(g_1, \ldots, g_m)$,
with $M = \mathcal{O}_{Y, y}^\wedge$, and with $N = S$
to conclude that the map is injective.
\end{proof}
\end{document}
